\section{A variable-shift relation for ordinary values}\label{sec:bridge}
We prove the ordinary-value deduction from the exact index model,
the compiled precision relation, the height theorem and $\mathsf{PC}$.
These are supplied by Sections~\ref{sec:model}, \ref{sec:compiler},
\ref{sec:height} and Theorem~\ref{hyp:PC}, respectively. Capital letters below denote five-coordinate point codes;
their lower-case indices are ordinary integers. The external variable
$t$ is any rational number.

\subsection{Approximation for every fixed shift}
Define $\Close(N,\mathcal K,t)$ by choosing $A,L,J$, linking all five
signed codes to precision inputs in the order $(n,\kappa,a,l,j)$, and
imposing $|t-r(a,j)|<\epsilon_l$ on those same coordinates. Put
\[
 B(n,\kappa)=H_6(S_{\rm ht})|\kappa F(n)|^{24}.
\]
The precision contract gives
\begin{align}
 \Close(\iota(n),\iota(\kappa),t)&\Longrightarrow
                  |t-n|<e^{-B(n,\kappa)},\label{eq:close-sound}\\
 \Close(\iota(n),\iota(\kappa),n)&\quad\text{for every }
                  n\notin\Cset,\ \kappa\ne0.\label{eq:close-complete}
\end{align}
The first follows from $|t-n|<2\epsilon_l<4^{-B}/2<e^{-B}$.
For the second, positive multiples of $P_0$ approach the identity
